% ACM conference paper. Use [acmsmall] or [acmlarge] for journals, and add
% `review` or `anonymous` for submission.
\documentclass[sigconf,nonacm]{acmart}

\usepackage{booktabs}

\begin{document}

\title{An Example Title}

\author{A. U. Thor}
\email{author@example.org}
\affiliation{%
  \institution{Your University}
  \city{City}
  \country{Country}}

\author{Second Author}
\email{second@example.org}
\affiliation{%
  \institution{Another Institute}
  \city{City}
  \country{Country}}

\begin{abstract}
State the problem, the approach and the main result in one paragraph.
\end{abstract}

\keywords{first keyword, second keyword, third keyword}

\maketitle

\section{Introduction}

Say what is known, what is missing, and what this paper adds. Cite earlier
work by key~\cite{einstein1905,knuth1984}.

\section{Method}

Describe what you did in enough detail to repeat it.

\begin{equation}
  \bar{x} = \frac{1}{n} \sum_{i=1}^{n} x_i
  \label{eq:mean}
\end{equation}

\begin{table}
  \caption{Example measurements}
  \label{tab:results}
  \begin{tabular}{lrr}
    \toprule
    Sample & Mean & Std.\ dev. \\
    \midrule
    A & 1.02 & 0.11 \\
    B & 0.97 & 0.08 \\
    \bottomrule
  \end{tabular}
\end{table}

\section{Results}

Report what you found. Table~\ref{tab:results} lists the measurements and
Equation~\ref{eq:mean} defines the mean.

\begin{acks}
Thank the people and grants that made the work possible.
\end{acks}

\bibliographystyle{ACM-Reference-Format}
\bibliography{references}

\end{document}
